%% ma: L12-B03-0018
%% lop: 12
%% chuong: 1
%% bai: 3
%% dang: 12.3.13
%% loai: TN4
%% muc: TH
%% dap_an: "C"
%% nguon: "HK1-DeOn_1_1.pdf (M10, Đề ôn tập học kì 1 Toán 12, Đề số 2), Phần I Câu 11"
%% kien_thuc: "Bài 3 (đường tiệm cận)"
%% trang_thai: chinh-thuc
%% ghi_chu: "Đã duyệt 29/10/2026."
\begin{ex}
Cho hàm số $y=\dfrac{x^2+1}{x^2+2x-3}$. Tổng số các đường tiệm cận của đồ thị hàm số đã cho là
\choice{$1$}{$2$}{\True $3$}{$4$}
\loigiai{$x^2+2x-3=0\Leftrightarrow x=1$ hoặc $x=-3$, tử số khác $0$ tại hai giá trị này nên có $2$ tiệm cận đứng $x=1$, $x=-3$. $\lim\limits_{x\to\pm\infty}y=1$ nên có tiệm cận ngang $y=1$. Tổng cộng $3$ đường.}
\end{ex}
